
\documentclass[11pt,oneside]{book}
\usepackage{../preamble}

\title{Mathematical Foundations of Idea Theory\\
\large Volume 1 of the IdeaTheory monograph}
\author{The IdeaTheory Project}
\date{\today}

\begin{document}
\frontmatter
\maketitle

\begin{abstract}
This volume lays the purely mathematical foundations of \emph{Idea Theory},
the formal framework developed throughout the six-volume IdeaTheory
monograph. An \emph{idea} is modelled as an element of a graded monoid
$(I,\circ,e,\deg)$ equipped with a bilinear \emph{resonance pairing}
$\rho : I \times I \to \mathbb{R}$, an \emph{Aufhebung decomposition}
$a \circ b = \mathsf{pres}(a,b) + \mathsf{neg}(a,b) + \mathsf{syn}(a,b)$,
and an \emph{emergence 2-cocycle} $\varepsilon$ measuring the failure of
composition to be strictly additive on a chosen valuation. From these data
we derive nine theorems whose statements and proofs are entirely
mathematical: (1) the composition theorem establishing the monoid
structure; (2) the bilinearity and non-degeneracy of the resonance
pairing; (3) the existence and uniqueness of the Aufhebung decomposition;
(4) the 2-cocycle identity for emergence; (5) the curvature theorem
relating asymmetry of $\rho$ to a closed 2-form on $I$; (6) the
conjugation theorem describing the action of invertible ideas by inner
automorphism; (7) the chain theorem on iterated composition and
transmission; (8) the structural-equivalence theorem characterising
isomorphism classes of idea algebras; and (9) the grading theorem
asserting that $\deg$ extends to a homomorphism into a totally ordered
abelian group.

The volume is deliberately self-contained and prerequisite-light: only
elementary linear algebra, basic group and monoid theory, and a working
familiarity with bilinear forms and 2-cocycles in the sense of group
cohomology are assumed. Every definition is given twice — once in
ordinary mathematical prose, once as it appears in the accompanying
Lean~4 formalisation under the namespace \texttt{IdeaTheory} — and every
theorem is stated with an explicit cross-reference to its machine-checked
counterpart. A dedicated \emph{From Informal Claims to Formal Statements}
table appears in each chapter, making the bridge between the
non-formalised literature on ideas, concepts, and emergence and the
fully verified Lean theorems entirely transparent.

The remaining five volumes of the monograph reinterpret this single
algebraic apparatus in five different domains: the social sciences
(Volume~2), the humanities (Volume~3), cognitive science and the
philosophy of mind (Volume~4), the philosophical problem of emergence
(Volume~5), and a collection of advanced and unifying applications
(Volume~6). The mathematics proved here is therefore the common spine
of the entire project: every applied claim made in Volumes~2--6 is, in
the end, an instance of one of the nine theorems established in this
volume. Readers who care only about the formal core may stop after
Volume~1; readers interested in interpretation will find that no further
mathematics is needed beyond what is proved here.
\end{abstract}

\tableofcontents

\chapter*{Preface to Volume~1}
\addcontentsline{toc}{chapter}{Preface to Volume~1}

\section*{Motivation}
Talk of ``ideas'' pervades the social sciences, the humanities,
cognitive science, and philosophy, yet the word has never had a
mathematically respectable referent. Authors as different as
Lovejoy~\cite{lovejoy1936}, Sperber~\cite{sperber1996}, Dennett~\cite{dennett1995},
and Deacon~\cite{deacon2012} use ``idea'' to denote something that can
be \emph{combined}, that can \emph{resonate} with or \emph{negate}
another idea, that can \emph{emerge} from simpler constituents, and
that can be \emph{transmitted} along chains of agents. Idea Theory
takes these four verbs at face value and asks: \emph{is there a single
algebraic object on which all four operations live, and in which the
familiar informal claims become provable theorems?} Volume~1 answers
yes, and proves the nine theorems that justify the answer.

\section*{Intended audience and prerequisites}
This volume is written for mathematicians, formal-methods researchers,
and mathematically inclined philosophers and social scientists. The
prerequisites are modest: undergraduate abstract algebra (groups,
monoids, rings), undergraduate linear algebra (bilinear forms, dual
spaces), and a passing acquaintance with the language of category
theory and group cohomology — chiefly the notion of a 2-cocycle. No
prior exposure to Lean~4 is required to read the prose; readers who
wish to inspect the formalisation will find a gentle introduction in
Chapter~1, and the project's \texttt{lakefile.lean} compiles against
a current Mathlib release.

\section*{Reading order}
The volume contains two chapters. Chapter~1 fixes notation,
introduces the data $(I,\circ,e,\deg,\rho,\mathsf{pres},
\mathsf{neg},\mathsf{syn},\varepsilon)$ of an \emph{idea algebra},
and proves Theorems~1--5 (composition, pairing, decomposition,
cocycle, curvature). Chapter~2 proves Theorems~6--9 (conjugation,
chains, equivalence, grading) and closes with a structure theorem
classifying finite-dimensional idea algebras up to isomorphism.
A reader in a hurry may read \S1.1--\S1.3 and \S2.4 and skip the
rest on first pass; a reader wishing to use Volume~1 as a reference
for any of Volumes~2--6 should read it linearly.

\section*{Relation to the other five volumes}
Volume~2 (\emph{Idea Theory and the Social Sciences}) interprets
$\circ$ as cultural transmission and $\rho$ as ideological
alignment. Volume~3 (\emph{Idea Theory and the Humanities}) reads
the Aufhebung decomposition as a model of dialectical reading and
intertextuality. Volume~4 (\emph{Cognitive Science and the
Philosophy of Mind}) identifies the curvature 2-form with a measure
of conceptual incongruity in mental representation. Volume~5
(\emph{Emergence: A Formal Theory}) is a sustained meditation on the
2-cocycle $\varepsilon$ as the precise mathematical content of the
philosophical notion of emergence. Volume~6 (\emph{Advanced
Applications and Unifying Perspectives}) collects applications that
do not fit cleanly into a single discipline. \emph{None of those
volumes prove any new mathematics}: every formal claim they invoke
is one of the nine theorems below.

\section*{Bridge to the Lean~4 formalisation}
Every definition, lemma, and theorem in this volume has a
machine-checked counterpart under the Lean namespace
\texttt{IdeaTheory}. The appendix to this volume lists every Lean
file relevant to Volume~1 with a one-paragraph description, and each
chapter contains a \emph{From Informal Claims to Formal Statements}
translation table mapping prose statements to their Lean names.

\mainmatter

\input{ch1}
\input{ch2}

\appendix

\chapter{Lean~4 Files Relevant to Volume~1}
\label{app:lean-vol1}

The following Lean~4 source files, all under the namespace
\texttt{IdeaTheory}, contain the machine-verified counterparts of the
definitions and theorems proved in this volume. Each is part of the
project's Lake build and depends only on a current Mathlib release.

\paragraph{\texttt{IdeaTheory/Foundations.lean}.}
Foundational declarations shared by every later file: the carrier type
\lstinline|Idea|, the composition operation \lstinline|comp|, the unit
\lstinline|e|, the grading map \lstinline|deg|, and basic notation.
Provides the \lstinline|IdeaAlgebra| typeclass that bundles the data of
an idea algebra together with the axioms used in Theorems~1--9.

\paragraph{\texttt{IdeaTheory/Composition.lean}.}
Proves Theorem~1 (\emph{idea composition monoid}): associativity, unit
laws, and cancellation properties of $\circ$. Docstrings cite the
informal composition claims in the literature on cultural transmission
and conceptual combination.

\paragraph{\texttt{IdeaTheory/Resonance.lean}.}
Proves Theorem~2 (\emph{resonance bilinear pairing}): bilinearity,
non-degeneracy, and the relationship between $\rho$ and the curvature
form. Bridges to the informal notion of ``ideological alignment''.

\paragraph{\texttt{IdeaTheory/Aufhebung.lean}.}
Proves Theorem~3 (\emph{Aufhebung decomposition}): existence and
uniqueness of the splitting
$a\circ b = \mathsf{pres}(a,b)+\mathsf{neg}(a,b)+\mathsf{syn}(a,b)$.
Cites Hegelian and post-Hegelian discussions of dialectical
sublation as the informal source.

\paragraph{\texttt{IdeaTheory/Emergence.lean}.}
Proves Theorem~4 (\emph{emergence 2-cocycle}): the 2-cocycle identity
$\delta\varepsilon = 0$ and a normalisation lemma. Cites the
philosophical literature on weak and strong emergence.

\paragraph{\texttt{IdeaTheory/Curvature.lean}.}
Proves Theorem~5 (\emph{meaning curvature form}): the antisymmetric
part of $\rho$ defines a closed 2-form whose cohomology class is
invariant under structural equivalence.

\paragraph{\texttt{IdeaTheory/Conjugation.lean}.}
Proves Theorem~6 (\emph{conjugation action}): invertible ideas act on
$I$ by inner automorphisms preserving $\rho$, $\mathsf{pres}$,
$\mathsf{neg}$, $\mathsf{syn}$, and $\varepsilon$.

\paragraph{\texttt{IdeaTheory/Chains.lean}.}
Proves Theorem~7 (\emph{transmission chain}): iterated composition is
controlled by a telescoping sum of cocycle values, yielding the
chain-rule formula used throughout Volumes~2 and~4.

\paragraph{\texttt{IdeaTheory/Equivalence.lean}.}
Proves Theorem~8 (\emph{structural equivalence}): two idea algebras
are isomorphic iff their pairings, decompositions, and cocycles agree
up to a graded monoid isomorphism.

\paragraph{\texttt{IdeaTheory/Grading.lean}.}
Proves Theorem~9 (\emph{graded idea algebra}): $\deg$ extends to a
homomorphism into a totally ordered abelian group, and the resulting
grading is compatible with $\rho$, the Aufhebung decomposition, and
$\varepsilon$.

\bibliographystyle{alpha}
\bibliography{../references}

\end{document}
=== END FILE ===
